
\chapter{相关技术与理论基础}

\section{LLM Agent技术}

\subsection{大语言模型与Agent范式}

大语言模型（Large Language Model，LLM）是基于Transformer架构\cite{Vaswani2017Attention}训练的大规模神经网络语言模型，通过在海量语料上进行预训练，在文本理解、代码生成、逻辑推理和知识问答等任务上表现出色。GPT-4\cite{OpenAI2024GPT4}、DeepSeek\cite{DeepSeekAI2024DeepSeekV3}等模型在多项自然语言处理上达到或超过人类水平。

Agent（智能体）在此基础上继续发展，将LLM由问答工具扩展为任务执行者。LLM Agent通常包含三个核心部分：感知模块接收环境数据和用户指令，决策模块推理出下一步动作，执行模块调用工具完成具体操作。Agent可以在没有人为干预的情况下自主执行多步骤的目标导向任务。

\subsection{ReAct框架}

ReAct（Reasoning and Acting）\cite{Yao2022ReAct}是Yao等人提出的一种将推理（Reasoning）与行动（Acting）交替进行的Agent执行框架。在ReAct框架中，Agent的每一步执行由三个要素组成：

\begin{itemize}
  \item Thought（推理）：Agent对当前状态进行分析，决定下一步应当采取的行动；
  \item Action（行动）：Agent调用一个外部工具执行具体操作；
  \item Observation（观察）：工具返回执行结果，Agent将其纳入上下文继续推理。
\end{itemize}

这三个要素组成一轮迭代过程，Agent 经过多轮Thought-Action-Observation循环逐步推进任务，直到产生Final Answer。与纯思维链推理方案相比，ReAct让Agent在每一步都能拿到工具的真实信息再做判断，减少了幻觉问题；与纯行动方案相比，显式的Thought步骤使决策过程可追溯。

\subsection{LangChain与LangGraph}

LangChain\cite{LangChain}是当前最广泛使用的LLM应用开发框架，它提供了对多种LLM的统一接口、工具调用抽象以及ReAct Agent的标准实现。LangChain通过create\_react\_agent函数创建Agent，将提示词模板、工具列表和LLM 名称作为参数接收，由AgentExecutor驱动执行循环。LangChain内置的ReActSingleInputOutputParser 能够从LLM的输出中解析Action和Action Input字段，得到工具调用指令。

LangGraph是LangChain生态中用于构建多Agent和复杂工作流的框架，是基于有向图的编程范式，支持定义节点（Agent或函数）、边（状态转移）和条件分支（动态路由）。LangGraph的核心抽象是StateGraph，开发者可以通过add\_node添加处理节点，通过add\_conditional\_edges定义基于状态的条件路由，从而构建灵活的多阶段执行流程。AutoPT的渗透测试状态机（PSM）即基于LangGraph实现，将Scan、VulnSelect、Inquire、Exploit、Check五个阶段建模为图中的节点，通过条件边控制状态转移。

\section{自动化渗透测试技术}

\subsection{渗透测试基本流程}

渗透测试（Penetration Testing）通过模拟攻击者行为来评估目标系统的安全性。一次完整的渗透测试通常包含以下阶段\cite{Liu2024AutoScan}\cite{Wu2023PenTest}：

（1）信息收集与侦察。使用端口扫描（如nmap）、服务指纹识别、目录遍历等手段，摸清目标系统的开放端口、运行服务、操作系统类型和Web应用框架。

（2）漏洞发现与分析。在信息收集的基础上，使用漏洞扫描工具（如xray、Nessus）对目标进行自动化扫描，或结合已知漏洞数据库（如NVD、ExploitDB）手动比对，识别目标系统可能存在的安全缺陷。

（3）漏洞利用。根据发现的漏洞，从公开渠道获取或编写PoC（Proof of Concept）代码，对目标系统实施攻击，尝试获取系统权限、读取敏感文件或执行任意命令。

（4）结果验证与报告。确认漏洞利用是否成功，收集利用证据（如打印/etc/passwd 文件内容、命令执行输出等），生成渗透测试报告。

传统渗透测试中，上述每个阶段都依赖安全人员的专业判断和手工操作。自动化渗透测试的目标是将其中可以程序化的部分交由系统完成，减少对人工的依赖。

\subsection{常用渗透测试工具}

本系统在不同阶段用到了以下工具。

nmap是渗透测试中最常见的网络扫描工具，覆盖端口发现、服务版本识别和操作系统指纹检测。典型用法是先用-sS -p-对目标做全端口SYN扫描，找到开放端口后再用-sV -sC识别各端口上运行的服务版本并执行默认NSE脚本。

xray是一款Web漏洞扫描器，可检测SQL注入、XSS、命令注入、文件读取等Web漏洞，有主动扫描和被动代理两种工作模式。本系统用它的主动扫描模式（xray ws --url）对目标HTTP服务做漏洞检测。

Playwright是微软开源的浏览器自动化框架，支持Chromium、Firefox和WebKit三种浏览器引擎。在渗透测试中，它主要用来模拟用户的浏览器操作，包括页面导航、表单填写、按钮点击等,处理那些需要Web界面交互才能完成的漏洞利用步骤。

\section{渗透测试状态机（PSM）}

\subsection{PSM的形式化定义}

AutoPT\cite{Wu2024AutoPT}用有限状态机（Finite State Machine，FSM）来描述渗透测试的执行流程，称之为渗透测试状态机（PSM）。其形式化定义是一个五元组：

PSM = (S, A, T, s₀, F)

S是状态集合，对应渗透测试中的各个阶段；A是动作集合，主要就是调用各种安全测试工具；T是状态转移函数（T: S×A→S），决定执行某个动作之后系统该跳到哪个状态；s₀是初始状态，F是终止状态集合。这套定义把一个本来连续的、难以把控的渗透测试过程拆成了离散的阶段，每个阶段该做什么、什么时候算结束都有了明确的界定。

\subsection{AutoPT的PSM状态设计}

AutoPT的PSM由五个状态串联而成，依次覆盖了信息收集、漏洞筛选、情报检索、漏洞利用和结果验证几个环节。Scan负责端口扫描和服务识别，输出攻击面信息；VulnSelect从扫描结果里筛选出值得尝试的漏洞；Inquire在正式利用之前去GitHub等外部渠道检索对应的PoC和利用方法；Exploit拿着这些信息执行具体的漏洞利用操作；Check验证利用是否成功，如果失败了，可以回退到Exploit换个方式再试，或者跳回VulnSelect换一个漏洞。具体状态转移关系见\cref{fig:psm_original}。

\begin{figure}[htbp]
\centering
\includegraphics[width=0.6\textwidth]{image6}
\caption{原版PSM状态转移图}
\label{fig:psm_original}
\end{figure}

整个状态机基于LangGraph实现，五个状态各自是图中的一个节点，节点之间的跳转由路由函数根据当前状态的输出动态决定。阅读AutoPT的源码可以发现，Check阶段的这个回退设计是它和纯线性流程最大的区别，Agent不会因为一次利用失败就直接退出，而是有机会调整方向重新来过。

\subsection{原始PSM的局限性}

虽然AutoPT通过PSM对渗透测试流程进行了结构化约束，但在实际复现几组实验后，我们发现它还有几处明显的缺陷。

（1）PSM流程控制机制存在多种问题。第一，状态粒度较粗，Scan状态把端口扫描和漏洞初筛混合在一起，一旦其中某个环节出问题，很难判断到底卡在哪一步。第二，状态转移条件过于固定。各状态之间基本按预定义的顺序推进，即使Scan阶段什么都没扫到，系统仍会构造一个占位语句传递给下游，后续的Inquire和Exploit阶段会跟着白跑一轮，造成不必要的Token消耗。第三，各Agent的提示词缺乏统一规范和角色区分，导致输出格式不稳定，Scan Agent有时候会越权去执行利用命令，Inquire Agent偶尔也会调用不属于它的工具。

（2）工具体系有缺口。Playwright工具包在原版本中没有表单填写和下拉框选择等功能，当遇到需要安装向导来配置环境的漏洞利用情况时，Agent就无法进行前置操作。ReAct解析器不能处理格式偏差，当使用思维链（Chains of Thought, CoT）模型进行分析时经常出现解析失败的情况。

（3）系统的可用性有待提升。AutoPT的交互完全基于命令行，用户需要通过编辑配置文件来设置测试参数，执行过程中只能看到文本输出，无法直观地监控测试进度和结果，也不方便对多个测试任务进行管理。

第三章将针对上述三个问题分别给出改进方案。

